% ===== ElegantBook customizations for 《Agentes de IA em Profundidade》 =====
% Accent color comes from ElegantBook's cyan theme: \structurecolor = RGB(31,186,190)

% ElegantBook loads Babel's Portuguese module when built with `lang=pt`.
% Document labels are customized below without reloading Babel with new options.

\AtBeginDocument{%
  \renewcommand{\contentsname}{Sumário}%
  \renewcommand{\chaptername}{Capítulo}%
  \renewcommand{\figurename}{Figura}%
  \renewcommand{\tablename}{Tabela}%
  \renewcommand{\listfigurename}{Lista de Figuras}%
  \renewcommand{\listtablename}{Lista de Tabelas}%
}

% ── Text & math fonts: Palatino look, matching the author's papers ──
\usepackage{unicode-math}
\setmainfont{texgyrepagella}[
  UprightFont = *-regular,
  BoldFont = *-bold,
  ItalicFont = *-italic,
  BoldItalicFont = *-bolditalic,
  Extension = .otf,
  Scale = 1.0]
\setsansfont{texgyreheros}[
  UprightFont = *-regular,
  BoldFont = *-bold,
  ItalicFont = *-italic,
  BoldItalicFont = *-bolditalic,
  Extension = .otf,
  Scale = 0.9]
\setmathfont{texgyrepagella-math}[Extension = .otf]

% ── Unicode glyph fallback (the text/mono fonts lack these) ──────
\usepackage{newunicodechar}
\IfFontExistsTF{Arial Unicode MS}{\newfontfamily\fallbackfont{Arial Unicode MS}}{\newfontfamily\fallbackfont{DejaVu Sans}}
\IfFontExistsTF{Heiti SC}{\newunicodechar{★}{{\fontspec{Heiti SC}★}}}{\newunicodechar{★}{{\fontspec{Noto Sans CJK SC}★}}}
\newunicodechar{✓}{{\fallbackfont ✓}}
\newunicodechar{✗}{{\fallbackfont ✗}}
\newunicodechar{₀}{{\fallbackfont ₀}}
\newunicodechar{₁}{{\fallbackfont ₁}}
\newunicodechar{₂}{{\fallbackfont ₂}}
\newunicodechar{₃}{{\fallbackfont ₃}}
\newunicodechar{₄}{{\fallbackfont ₄}}
\newunicodechar{₅}{{\fallbackfont ₅}}
\newunicodechar{₆}{{\fallbackfont ₆}}
\newunicodechar{₇}{{\fallbackfont ₇}}
\newunicodechar{₈}{{\fallbackfont ₈}}
\newunicodechar{₉}{{\fallbackfont ₉}}
\newunicodechar{ⓐ}{{\fallbackfont ⓐ}}
\newunicodechar{ⓑ}{{\fallbackfont ⓑ}}
\newunicodechar{ⓒ}{{\fallbackfont ⓒ}}
\newunicodechar{ⓓ}{{\fallbackfont ⓓ}}
\newunicodechar{ⓔ}{{\fallbackfont ⓔ}}
\newunicodechar{①}{{\fallbackfont ①}}
\newunicodechar{②}{{\fallbackfont ②}}
\newunicodechar{③}{{\fallbackfont ③}}
\newunicodechar{④}{{\fallbackfont ④}}
\newunicodechar{⑤}{{\fallbackfont ⑤}}
\newunicodechar{⑥}{{\fallbackfont ⑥}}
\newunicodechar{⑦}{{\fallbackfont ⑦}}
\newunicodechar{⑧}{{\fallbackfont ⑧}}
\newunicodechar{⑨}{{\fallbackfont ⑨}}
\newunicodechar{⑩}{{\fallbackfont ⑩}}
% Math / symbols absent from the Latin text font
\newunicodechar{≈}{{\fallbackfont ≈}}
\newunicodechar{♠}{{\fallbackfont ♠}}
\newunicodechar{♣}{{\fallbackfont ♣}}
\newunicodechar{♥}{{\fallbackfont ♥}}
\newunicodechar{♦}{{\fallbackfont ♦}}
% Greek letters used inline in body text
\newunicodechar{α}{{\fallbackfont α}}
\newunicodechar{β}{{\fallbackfont β}}
\newunicodechar{γ}{{\fallbackfont γ}}
\newunicodechar{δ}{{\fallbackfont δ}}
\newunicodechar{ε}{{\fallbackfont ε}}
\newunicodechar{η}{{\fallbackfont η}}
\newunicodechar{θ}{{\fallbackfont θ}}
\newunicodechar{λ}{{\fallbackfont λ}}
\newunicodechar{μ}{{\fallbackfont μ}}
\newunicodechar{π}{{\fallbackfont π}}
\newunicodechar{ρ}{{\fallbackfont ρ}}
\newunicodechar{σ}{{\fallbackfont σ}}
\newunicodechar{τ}{{\fallbackfont τ}}
\newunicodechar{φ}{{\fallbackfont φ}}
\newunicodechar{ω}{{\fallbackfont ω}}

% ── Residual CJK glyphs ──────────────────────────────────────────
\usepackage{xeCJK}
\IfFontExistsTF{Heiti SC}{%
  \setCJKmainfont{Heiti SC}%
  \setCJKmonofont{Heiti SC}%
}{%
  \setCJKmainfont{Noto Sans CJK SC}%
  \setCJKmonofont{Noto Sans Mono CJK SC}%
}

% ── Code/monospace font (covers box-drawing glyphs used in code) ──
\IfFontExistsTF{Menlo}{\setmonofont[Scale=0.86]{Menlo}}{\setmonofont[Scale=0.86]{DejaVu Sans Mono}}

% ── tcolorbox (ElegantBook already loads it; just add libraries) ──
\tcbuselibrary{skins,breakable}

% ── TikZ (for the full-bleed image cover in cover.tex) ───────────
\usepackage{tikz}
\usetikzlibrary{fadings}
\tikzfading[name=titlefade, top color=transparent!100, bottom color=transparent!0]

% Theme accent color: override ElegantBook's bright cyan with a deep navy
\definecolor{structurecolor}{RGB}{30,58,107}

% Accent shade for code/figure chrome
\colorlet{accent}{structurecolor}

% Internal cross-reference links
\newcommand{\crossreflink}[2]{\hyperref[#1]{\textcolor{structurecolor}{#2}}}
\colorlet{accentbg}{structurecolor!6}
\colorlet{accentrule}{structurecolor!35}

% ── Wrap long lines in code blocks ───────────────────────────────
\usepackage{fvextra}
\fvset{breaklines=true,breakanywhere=true,%
  breaksymbolleft={\footnotesize\textcolor{accentrule}{$\hookrightarrow$}\,}}
\DefineVerbatimEnvironment{Highlighting}{Verbatim}{commandchars=\\\{\},breaklines,breakanywhere}

% ── Code blocks: subtle O'Reilly-style framed box ────────────────
\newtcolorbox{codebox}{%
  breakable, enhanced,
  colback=black!3, colframe=accentrule,
  boxrule=0.5pt, leftrule=2.5pt,
  left=0.7em, right=0.6em, top=0.4em, bottom=0.4em,
  arc=0.6mm,
  before skip=0.7em, after skip=0.7em
}
\renewenvironment{Shaded}{\begin{codebox}}{\end{codebox}}
\renewenvironment{verbatim}%
  {\VerbatimEnvironment\begin{codebox}\begin{Verbatim}}%
  {\end{Verbatim}\end{codebox}}

% ── Experiment & question callouts (used by experiment_box.lua) ───
\newtcolorbox{experimentbox}{
  breakable, enhanced,
  colback=accentbg, colframe=accent,
  boxrule=0.6pt, left=0.8em, right=0.8em, top=0.6em, bottom=0.6em,
  arc=1mm, before skip=0.9em, after skip=0.9em
}

\newtcolorbox{questionbox}{
  breakable, enhanced,
  colback=accentbg, colframe=accent,
  boxrule=0.8pt, left=0.9em, right=0.9em, top=0.9em, bottom=0.7em,
  arc=1.2mm, before skip=0.9em, after skip=0.9em,
  title={\sffamily\bfseries Questões para reflexão},
  fonttitle=\normalsize\color{white},
  coltitle=white,
  attach boxed title to top left={xshift=1em, yshift=-\tcboxedtitleheight/2},
  boxed title style={colback=accent, colframe=accent, arc=0.8mm, outer arc=0.8mm}
}

% ── Figures: force [H] so they can sit inside callout/quote boxes ─
\usepackage{float}
\let\origfigure\figure
\let\origendfigure\endfigure
\renewenvironment{figure}[1][htbp]{\origfigure[H]}{\origendfigure}
\setkeys{Gin}{width=\linewidth,height=0.38\textheight,keepaspectratio}

% ── Figure captions: manual numbering in text, so suppress label ──
\usepackage{caption}
\captionsetup{font={normalsize,sf,bf}, labelformat=empty, skip=4pt, justification=centering}

% ── Blockquotes: teal left bar (experiments render as quotes) ────
\renewenvironment{quote}{%
  \begin{tcolorbox}[
    blanker, breakable,
    borderline west={2.5pt}{0pt}{accent},
    colback=accent!4,
    left=1em, right=0.6em, top=0.5em, bottom=0.5em,
    before skip=0.6em, after skip=0.6em
  ]%
}{%
  \end{tcolorbox}%
}

% ── Inline code: keep it tight ───────────────────────────────────
\setlength{\parindent}{2em}

% ElegantBook empties the plain page style used on chapter-opening pages.
\fancypagestyle{plain}{%
  \fancyhf{}%
  \fancyfoot[C]{\color{structurecolor}\small\thepage}%
  \renewcommand{\headrulewidth}{0pt}%
  \renewcommand{\footrulewidth}{0pt}%
}
